\section{Follow the dependency through one quarter-round}
A quarter-round updates four words. Each update uses the values produced by the
preceding updates. The calculation therefore has an order even when its words occupy
separate registers. Load the words into \texttt{t0}, \texttt{t1}, \texttt{t2} and
\texttt{t3}, representing \(a,b,c,d\). The first group adds \(b\) into \(a\), XORs
that new \(a\) into \(d\), then rotates the new \(d\) left by sixteen.
\Needspace{8\baselineskip}
\MARSCodeLabel{The first dependent group changes a, then d.}
\begin{minted}[samepage=true]{mars-idea}
addu $t0, $t0, $t1
xor $t3, $t3, $t0
rotr $t3, $t3, 16
\end{minted}

The second group adds the new \(d\) into \(c\), XORs the new \(c\) into \(b\), and
rotates \(b\) left by twelve. The third and fourth groups repeat the two pairings,
with left rotations by eight and seven. In the listing, the right-rotate immediates
are the complementary amounts, 24 and 25.
\Needspace{14\baselineskip}
\MARSCodeLabel{Complete the quarter-round using the updated words.}
\begin{minted}[samepage=true]{mars-idea}
addu $t2, $t2, $t3
xor $t1, $t1, $t2
rotr $t1, $t1, 20
addu $t0, $t0, $t1
xor $t3, $t3, $t0
rotr $t3, $t3, 24
addu $t2, $t2, $t3
xor $t1, $t1, $t2
rotr $t1, $t1, 25
\end{minted}

Load all four source words before making any stores, then store the four final words
through their corresponding argument pointers. This leaf function needs no saved
working registers and makes no nested call. Returning through \texttt{ra} preserves
its full address. The disjoint-word contract keeps the four destinations unambiguous.

For the first numerical check, use \(a=\mathtt{11111111}\),
\(b=\mathtt{01020304}\), \(c=\mathtt{9b8d6f43}\) and
\(d=\mathtt{01234567}\). These are the quarter-round inputs in RFC 8439,
section 2.1.1. The following intermediate table is calculated by the independent
word model. Every entry is hexadecimal.

\begin{table}[htbp]
\centering\caption{Working words after each add, XOR and rotation group.}
\tagpdfsetup{table/header-rows={1}}
\begin{tabular}{rllll}\toprule
Group & \(a\) & \(b\) & \(c\) & \(d\)\\\midrule
1 & \texttt{12131415} & \texttt{01020304} & \texttt{9b8d6f43} & \texttt{51721330}\\
2 & \texttt{12131415} & \texttt{d8177edf} & \texttt{ecff8273} & \texttt{51721330}\\
3 & \texttt{ea2a92f4} & \texttt{d8177edf} & \texttt{ecff8273} & \texttt{5881c4bb}\\
4 & \texttt{ea2a92f4} & \texttt{cb1cf8ce} & \texttt{4581472e} & \texttt{5881c4bb}\\
\bottomrule\end{tabular}
\end{table}

In group one, XOR produces \texttt{13305172} before rotation. Swapping its two
16-bit halves gives \texttt{51721330}. In group two, the pre-rotation value is
\texttt{edfd8177}; a left rotation by twelve gives \texttt{d8177edf}. Stop after
each group and inspect only the two words it changes. If the first changed word
agrees and the second does not, inspect the XOR operands and rotate amount before
altering the addition.
